\section{Randomized trials, local privacy, and welfare}

Throughout, \leanref{sym:n}{\ensuremath{n\ge2}} denotes the participant count and \leanref{sym:d}{\ensuremath{d\ge2}} the stratum count. Stratum indices range over \leanref{sym:[d]}{\ensuremath{[d]=\{1,\ldots,d\}}}, and participant indices range from \(1\) to \(n\). Generic positive numerical constants are denoted by \leanref{sym:c}{\ensuremath{c}} and \leanref{sym:C_0}{\ensuremath{C_0}}; their values may differ across comparisons. For a finite vector, the \(\ell_1\) norm is the sum of its absolute coordinates. The notation \(\asymp\) denotes comparison within fixed positive multiplicative constants, with its sequence convention specified in \cref{obj:synth_3}. We first describe the trial population and welfare target, then the private observation scheme and the statistical criteria.

\subsection{Trial population and welfare}

Each participant has a stratum, a treatment assignment, an observed outcome, and two potential outcomes. We first state the public uniform-stratum condition; the complete record space and observed-record marginal are collected explicitly below.

\begin{assumptionv}{ass:uniform-covariates}[Uniform participant strata]
The participant stratum \(X\) is uniformly distributed on \([d]\):
\[
P(X=j)=\frac{1}{d}\qquad(j\in[d]).
\]
\end{assumptionv}

Thus the stratum \leanref{sym:X}{\ensuremath{X}}, assignment \leanref{sym:A}{\ensuremath{A}}, and observed outcome \leanref{sym:Y}{\ensuremath{Y}} form the observed record \leanref{sym:O}{\ensuremath{O}} in the alphabet \leanref{sym:\mathcal O}{\ensuremath{\mathcal O}}. The potential outcomes \leanref{sym:Y^0}{\ensuremath{Y^0}} and \leanref{sym:Y^1}{\ensuremath{Y^1}} enter the full-data law, while \leanref{sym:Y^{\mathrm{pot}}}{\ensuremath{Y^{\mathrm{pot}}_A}} denotes the outcome selected by assignment. The observed-record distribution is \leanref{sym:P_O}{\ensuremath{P_O}}, as defined in \cref{obj:synth_1}.

For laws assigning positive probability to every stratum, define the arm-specific conditional mean \leanref{sym:\mu}{\ensuremath{\mu_{aj}(P)}} and treatment contrast \leanref{sym:\tau}{\ensuremath{\tau_j(P)}} by
\[
\mu_{aj}(P)=\mathbb E_P[Y^a\mid X=j],
\qquad
\tau_j(P)=\mu_{1j}(P)-\mu_{0j}(P),
\qquad
\tau(P)=(\tau_1(P),\ldots,\tau_d(P)).
\]
Here the arm index \leanref{sym:a}{\ensuremath{a\in\{0,1\}}} distinguishes control from treatment, and \leanref{sym:j}{\ensuremath{j\in[d]}} indexes the stratum. A cellwise treatment rule selects an arm in each stratum. Under uniform strata, its first-best population welfare is obtained by choosing the arm with the larger conditional mean. We write the first-best value \leanref{sym:V}{\ensuremath{V(P)}} and the observed-outcome baseline \leanref{sym:B}{\ensuremath{B(P)}} as
\[
V(P)=\frac1d\sum_{j=1}^d\max\{\mu_{0j}(P),\mu_{1j}(P)\},
\qquad
B(P)=\mathbb E_P[Y].
\]
This target evaluates the welfare attainable by an optimal stratum-specific treatment choice.

The remaining single-participant restrictions specify randomized assignment, the relation between observed and potential outcomes, and the admissible conditional means.

\begin{assumptionv}{ass:fair-randomization}[Fair conditional treatment assignment]
The binary treatment assignment \(A\), conditional on the participant stratum \(X\) and potential outcomes \(Y^0,Y^1\), satisfies
\[
\mathcal L_P(A\mid X,Y^0,Y^1)=\operatorname{Bernoulli}(1/2).
\]
\end{assumptionv}

Uniform categorical sampling fixes the stratum probabilities at a public reference distribution, specializing the multinomial sampling framework used for distributional functionals \citep{JiaoHanWeissman2018L1}. Equal population weights make each cell's contribution to welfare explicit, while the realized participant stratum remains part of the private record.

\begin{assumptionv}{ass:consistency}[Observed outcome consistency]
The observed outcome \(Y\) equals \(Y^{\mathrm{pot}}_A\), the potential outcome selected by the realized assignment, almost surely:
\[
P(Y=Y^{\mathrm{pot}}_A)=1.
\]
\end{assumptionv}

Known randomized treatment assignment is maintained here in its fair-assignment form \citep{OhnishiAwan2025Randomized}. Conditioning on both potential outcomes ensures that assignment supplies experimental variation within each stratum and gives the two arms equal weight in the baseline mean.

\begin{assumptionv}{ass:interior-means}[Interior potential-outcome means]
The conditional potential-outcome means \(\mu_{aj}(P)\) satisfy
\[
\mu_{aj}(P)\in[1/4,3/4]
\qquad(a\in\{0,1\},\ j\in[d]).
\]
\end{assumptionv}

Potential-outcome consistency connects the recorded response to the outcome under the assigned arm, giving the observed data their treatment-welfare interpretation \citep{LuedtkeVanDerLaan2016OptimalValue}. Together with fair assignment, it identifies each conditional arm mean from the corresponding observed treatment group.

\begin{assumptionv}{ass:iid-people}[Independent identically distributed participants]
The vector of observed participant records \(\mathbf O\) has the product law
\[
\mathcal L_P(\mathbf O)=P_O^{\otimes n},
\]
where \(P_O\) is the observed-record marginal law induced by \(P\).
\end{assumptionv}

Fixed interior Bernoulli means strengthen the bounded-outcome framework for optimal treatment value \citep{LuedtkeVanDerLaan2016OptimalValue}. The maintained interval keeps every arm mean away from the outcome boundaries, places treatment contrasts in \([-1/2,1/2]\), and places first-best welfare in \([1/4,3/4]\).

We now define the complete record space before collecting the four single-participant restrictions into the causal class.

\begin{definitionv}{synth_1}[Full-data causal law space]
For an integer \(d\ge2\), write \([d]=\{1,\ldots,d\}\). Define \(\mathcal P_d\) as the class of all probability laws on the finite record space
\[
[d]\times\{0,1\}^4,
\]
equipped with its discrete sigma-algebra. The coordinate variables are ordered as \((X,A,Y,Y^0,Y^1)\): \(X\in[d]\) is the covariate stratum, \(A\in\{0,1\}\) is the treatment assignment, \(Y\in\{0,1\}\) is the observed outcome, and \(Y^0,Y^1\in\{0,1\}\) are the potential outcomes under control and treatment, respectively. Write \(Y^{\mathrm{pot}}_0=Y^0\) and \(Y^{\mathrm{pot}}_1=Y^1\). For \(P\in\mathcal P_d\), the observed record is \(O=(X,A,Y)\in\mathcal O=[d]\times\{0,1\}^2\), and \(P_O\) denotes its marginal law under \(P\).
\end{definitionv}

The causal class \leanref{sym:\mathcal V_d}{\ensuremath{\mathcal V_d}} in \cref{obj:def:causal-model} leaves the dependence between the two potential outcomes unrestricted subject to these conditions. Under this class, fair assignment and consistency give
\[
B(P)=\frac1{2d}\sum_{j=1}^d
\bigl(\mu_{0j}(P)+\mu_{1j}(P)\bigr).
\]
For a contrast vector, let the average absolute contrast \leanref{sym:F}{\ensuremath{F(\theta)}} be defined by
\[
F(\theta)=\frac1d\sum_{j=1}^d|\theta_j|.
\]
The identity \(\max\{u,v\}=(u+v+|u-v|)/2\) then yields
\[
V(P)=B(P)+\frac12F(\tau(P)).
\]
The baseline records welfare under fair assignment; the absolute-contrast component records the gain from selecting the better arm separately in each stratum.

The following definition collects the four single-participant restrictions. Sampling is specified separately through the iid condition above, with observed sample \leanref{sym:\mathbf O}{\ensuremath{\mathbf O=(O_1,\ldots,O_n)}}.

\begin{definitionv}{def:causal-model}[Causal law class $\mathcal V_d$]
\(\mathcal V_d\) is the class of full-data probability laws \(P\in\mathcal P_d\) satisfying all four conditions:
\begin{itemize}
\item \textbf{(Uniform strata.)} \Cref{obj:ass:uniform-covariates}.
\item \textbf{(Fair assignment.)} \Cref{obj:ass:fair-randomization}.
\item \textbf{(Consistency.)} \Cref{obj:ass:consistency}.
\item \textbf{(Interior means.)} \Cref{obj:ass:interior-means}.
\end{itemize}
\end{definitionv}

Independent identically distributed one-record participants provide the standard sampling structure for local-private minimax analysis \citep{DuchiJordanWainwright2018Minimax}. Every participant contributes one draw from the same trial population, and the collection protocol acts on these separate records.

\subsection{Private collection and statistical criteria}

Fix the privacy budget \leanref{sym:\varepsilon}{\ensuremath{0<\varepsilon\le1}}. Collection may use a shared public seed \leanref{sym:R}{\ensuremath{R}} taking values in a standard-Borel seed space \leanref{sym:\mathcal R}{\ensuremath{\mathcal R}}, with probability law \leanref{sym:\lambda}{\ensuremath{\lambda}}. Participant \leanref{sym:i}{\ensuremath{i}} sends one message in a standard-Borel message space \leanref{sym:\mathcal Z}{\ensuremath{\mathcal Z_i}}. The analyst also has a uniform randomizer \leanref{sym:U}{\ensuremath{U}} for randomized statistical decisions. We first state the independence condition on public and analyst randomness.

\begin{assumptionv}{ass:independent-randomness}[Independent public and analyst randomness]
The shared public seed \(R\), with probability law \(\lambda\), and the analyst randomizer \(U\) have conditional joint law
\[
\mathcal L(R,U\mid\mathbf O)
=\lambda\otimes\operatorname{Uniform}[0,1],
\]
where \(\mathbf O\) is the vector of observed participant records.
\end{assumptionv}

The history space \leanref{sym:\mathcal H}{\ensuremath{\mathcal H_i}} in \cref{obj:synth_2} contains the messages preceding participant \(i\). A history \leanref{sym:\eta}{\ensuremath{\eta}} and a seed realization \leanref{sym:r}{\ensuremath{r}} therefore index the participant's conditional message distribution \leanref{sym:Q}{\ensuremath{Q_i(\cdot\mid o,\eta,r)}}, where \leanref{sym:o}{\ensuremath{o}} denotes an observed-record input. Sequential collection generates
\[
Z_i\mid \mathbf O,R,Z_1,\ldots,Z_{i-1}
\sim Q_i(\cdot\mid O_i,(Z_1,\ldots,Z_{i-1}),R).
\]
The private transcript is the message vector \leanref{sym:\mathbf Z}{\ensuremath{\mathbf Z=(Z_1,\ldots,Z_n)}}; the analyst observes this vector and the public seed. The next condition imposes privacy on every history-indexed message kernel.

\begin{assumptionv}{ass:full-record-privacy}[Full-record local privacy]
For every participant index \(i\), observed records \(o,o'\), history \(\eta\), public seed \(r\), and measurable message set \(E\) in their declared domains,
\[
Q_i(E\mid o,\eta,r)\le e^\varepsilon Q_i(E\mid o',\eta,r).
\]
\end{assumptionv}

Independent public coins and independent randomized decision rules allow coordination and randomized analysis within the observation scheme \citep{AcharyaCanonneFreitagTyagi2019TestWithoutTrust}. The product conditional law makes their joint distribution the same for every observed sample.

Privacy compares the message distributions induced by any two complete observed records, including their strata, assignments, and outcomes. Let \leanref{sym:o'}{\ensuremath{o'}} denote the second record and \leanref{sym:E}{\ensuremath{E}} a measurable message event. The following definition makes the history domain used in that inequality explicit.

\begin{definitionv}{synth_2}[Participant message histories]
For a protocol with \(n\) participants, let \((\mathcal Z_i)_{i=1}^n\) be its standard-Borel participant message spaces. For each \(i\in\{1,\ldots,n\}\), define the history space
\[
\mathcal H_i=\prod_{\ell=1}^{i-1}\mathcal Z_\ell
\]
with the product sigma-algebra; the empty product \(\mathcal H_1\) is a singleton. A history \(\eta\in\mathcal H_i\) records the messages of participants \(1,\ldots,i-1\). Thus, with public-seed space \(\mathcal R\), the participant kernel \(Q_i\) is a Markov kernel from \(\mathcal O\times\mathcal H_i\times\mathcal R\) to \(\mathcal Z_i\), and the realized history preceding participant \(i\) is \((Z_1,\ldots,Z_{i-1})\).
\end{definitionv}

Pure sequential local differential privacy controls the change in each participant's message distribution when the complete input record changes \citep{DuchiJordanWainwright2018Minimax}. Requiring the bound at every seed and history applies the same protection to every conditional collection rule, including histories that arise under other population laws.

The sequential class permits later collection rules to depend on earlier messages. Its noninteractive subclass imposes a further condition, with \leanref{sym:\eta\prime}{\ensuremath{\eta'}} denoting a second history.

\begin{assumptionv}{ass:noninteraction}[Independence from message history]
For every participant index \(i\), observed record \(o\), public seed \(r\), and histories \(\eta,\eta'\in\mathcal H_i\),
\[
Q_i(\cdot\mid o,\eta,r)=Q_i(\cdot\mid o,\eta',r).
\]
\end{assumptionv}

Noninteractive channels use rules fixed independently of the preceding messages, while permitting participant-specific rules and shared public coins \citep{AcharyaCanonneFreitagTyagi2019TestWithoutTrust}. Thus coordination through the independent seed remains available within this subclass.

We collect these requirements into two protocol classes, both of which give each participant one message.

\begin{definitionv}{def:sequential-class}[Sequential protocol class $\mathfrak Q_{\mathrm{SI}}(n,d,\varepsilon)$]
\(\mathfrak Q_{\mathrm{SI}}(n,d,\varepsilon)\) is the class of sequences
\[
Q=(Q_i)_{i=1}^n
\]
of the declared Markov kernels satisfying \cref{obj:ass:full-record-privacy}.
\end{definitionv}

The one-message sequential class \leanref{sym:\mathfrak Q_{\mathrm{SI}}}{\ensuremath{\mathfrak Q_{\mathrm{SI}}(n,d,\varepsilon)}} in \cref{obj:def:sequential-class} allows the conditional message distribution to respond to earlier releases. The next definition restricts this dependence.

\begin{definitionv}{def:noninteractive-class}[Noninteractive protocol class $\mathfrak Q_{\mathrm{NI}}(n,d,\varepsilon)$]
\(\mathfrak Q_{\mathrm{NI}}(n,d,\varepsilon)\) is the subclass of \(\mathfrak Q_{\mathrm{SI}}(n,d,\varepsilon)\) satisfying both \cref{obj:ass:full-record-privacy,obj:ass:noninteraction}.
\end{definitionv}

The noninteractive class \leanref{sym:\mathfrak Q_{\mathrm{NI}}}{\ensuremath{\mathfrak Q_{\mathrm{NI}}(n,d,\varepsilon)}} in \cref{obj:def:noninteractive-class} is contained in the sequential class. Let the class index \leanref{sym:\mathcal C}{\ensuremath{\mathcal C\in\{\mathrm{NI},\mathrm{SI}\}}} select the admissible protocol class \leanref{sym:\mathfrak Q}{\ensuremath{\mathfrak Q_{\mathcal C}(n,d,\varepsilon)}}. A welfare estimate \leanref{sym:T}{\ensuremath{T(\mathbf Z,R,U)}} is a Borel real-valued function of the messages, public seed, and analyst randomizer. Expectations and probabilities below include participant sampling, message generation, and the independent randomizers. Our precision criterion optimizes over both the collection protocol and the estimate.

\begin{definitionv}{def:risk}[Minimax welfare risk $\mathfrak R_{\mathcal C}(n,d,\varepsilon)$]
\[
\mathfrak R_{\mathcal C}(n,d,\varepsilon)
=
\inf_{Q\in\mathfrak Q_{\mathcal C}(n,d,\varepsilon)}
\inf_T
\sup_{P\in\mathcal V_d}
\mathbb E_{P,Q}
\left[
\bigl(T(\mathbf Z,R,U)-V(P)\bigr)^2
\right].
\]
Here \(\mathcal C\in\{\mathrm{NI},\mathrm{SI}\}\) selects the corresponding declared protocol class \(\mathfrak Q_{\mathcal C}\), and the second infimum ranges over Borel real-valued welfare estimates \(T(\mathbf Z,R,U)\) based on the private message vector \(\mathbf Z\), the public seed \(R\), and the independent uniform analyst randomizer \(U\). The target \(V(P)\) is the average of the larger arm mean in each stratum under \(P\).
\end{definitionv}

The minimax squared-error risk \leanref{sym:\mathfrak R}{\ensuremath{\mathfrak R_{\mathcal C}(n,d,\varepsilon)}} in \cref{obj:def:risk} measures the best worst-law precision available within the selected protocol class. Its population supremum permits arbitrary treatment-effect patterns satisfying the causal restrictions.

For inference, a welfare interval \leanref{sym:I}{\ensuremath{I(\mathbf Z,R,U)}} is connected and closed, with ordered Borel endpoints in the welfare range \([1/4,3/4]\). Its length is the upper endpoint minus the lower endpoint. Global honesty requires coverage of at least \(0.90\) at every law in the causal class; the corresponding criterion minimizes worst-law expected length.

\begin{definitionv}{def:honest-length}[Honest interval length $\mathfrak H_{\mathcal C}(n,d,\varepsilon)$]
The minimax expected length of a globally honest welfare interval is
\[
\mathfrak H_{\mathcal C}(n,d,\varepsilon)
=
\inf_{Q\in\mathfrak Q_{\mathcal C}(n,d,\varepsilon)}
\ \inf_{\substack{I:\ \inf_{P\in\mathcal V_d}
\Pr_{P,Q}\{V(P)\in I(\mathbf Z,R,U)\}\ge0.90}}
\ \sup_{P\in\mathcal V_d}
\mathbb E_{P,Q}\!\left[
\operatorname{length} I(\mathbf Z,R,U)
\right].
\]
The second infimum ranges over connected closed intervals whose ordered endpoints are Borel functions of the private messages \(\mathbf Z\), public seed \(R\), and analyst randomizer \(U\), and take values in \([1/4,3/4]\).
\end{definitionv}

The honest-length criterion \leanref{sym:\mathfrak H}{\ensuremath{\mathfrak H_{\mathcal C}(n,d,\varepsilon)}} in \cref{obj:def:honest-length} couples interval precision to uniform coverage over the same population class used for squared-error risk. Connectedness makes the reported uncertainty a single welfare range, and the endpoint restriction respects the range implied by the interior means.

\subsection{The signed experiment}

The welfare decomposition separates a population mean from an absolute-value functional of treatment contrasts. To connect the latter to a finite-alphabet experiment, define the signed treatment-outcome statistic \leanref{sym:S}{\ensuremath{S}} by
\[
S=(2A-1)(2Y-1).
\]
The signed input consists of the stratum and this sign, taking values in the paired alphabet \leanref{sym:\mathcal A_d}{\ensuremath{\mathcal A_d=[d]\times\{-1,1\}}}. For a contrast parameter \leanref{sym:\theta}{\ensuremath{\theta\in[-1/2,1/2]^d}}, define the paired law \leanref{sym:P_\theta}{\ensuremath{P_\theta}} and the uniform reference law \leanref{sym:\mathsf U_{2d}}{\ensuremath{\mathsf U_{2d}}} by
\[
P_\theta(j,s)=\frac{1+s\theta_j}{2d},
\qquad
\mathsf U_{2d}(j,s)=\frac1{2d}.
\]
Here \leanref{sym:s}{\ensuremath{s\in\{-1,1\}}} is the sign coordinate. Total variation on this finite alphabet is one half the sum of the absolute differences in probability masses, so
\[
\operatorname{TV}(P_\theta,\mathsf U_{2d})
=\frac1{2d}\sum_{j=1}^d|\theta_j|
=\frac12F(\theta).
\]
These definitions connect the absolute-contrast welfare component to distance from a public reference distribution.

A symmetric causal family embeds every allowed contrast vector into the trial model while keeping the baseline fixed. Its construction specifies the dependence of the potential outcomes as well as their means.

\begin{definitionv}{synth_3}[Private resource rates and their regimes]
For \(n,d\in\{2,3,\ldots\}\) and \(0<\varepsilon\le1\), put \(t=n\varepsilon^2\) and \(L=\log(ed)\). Define
\[
\rho(t,d)=
\begin{cases}
\dfrac{d^2}{tL},&t\ge d^2L,\\[4pt]
\displaystyle\min\left\{1,
\left[\dfrac{\log(e+d^2L/t)}{L}\right]^2\right\},
&0<t<d^2L,
\end{cases}
\qquad
r_{\mathcal C}(n,d,\varepsilon)=\rho(t,d),
\qquad
h_{\mathcal C}(n,d,\varepsilon)=\sqrt{\rho(t,d)}
\]
for each \(\mathcal C\in\{\mathrm{NI},\mathrm{SI}\}\).

For positive numerical sequences \(f_\ell,g_\ell\), write \(f_\ell\asymp g_\ell\) when there exist constants \(0<c\le C_0<\infty\) and an integer \(\ell_0\) such that
\[
c g_\ell\le f_\ell\le C_0g_\ell
\qquad(\ell\ge\ell_0).
\]
We say that these rates satisfy the resource conclusions when all the following properties hold.

For every integer \(d\ge2\) and every \(0<t<d^2\log(ed)\),
\[
\rho(t,d)=1
\quad\Longleftrightarrow\quad
t\le\frac{d^2\log(ed)}{ed-e}.
\]
Thus the saturated regime ends at \(d^2\log(ed)/(ed-e)\), the intermediate regime occupies
\[
\frac{d^2\log(ed)}{ed-e}<t<d^2\log(ed),
\]
and the dense regime starts at \(t=d^2\log(ed)\).

For every sequence \((n_\ell,d_\ell,\varepsilon_\ell)_{\ell\ge1}\) satisfying
\[
n_\ell,d_\ell\in\{2,3,\ldots\},
\qquad 0<\varepsilon_\ell\le1,
\]
put \(t_\ell=n_\ell\varepsilon_\ell^2\) and \(L_\ell=\log(ed_\ell)\). Then
\[
\rho(t_\ell,d_\ell)\longrightarrow0
\quad\Longleftrightarrow\quad
t_\ell\longrightarrow\infty
\ \text{and}\
\frac{\log(1+d_\ell^2/t_\ell)}{L_\ell}\longrightarrow0,
\]
and
\[
\rho(t_\ell,d_\ell)\asymp t_\ell^{-1}
\quad\Longleftrightarrow\quad
\left(\exists a>0,\ \exists\ell_0,\ \forall\ell\ge\ell_0:
t_\ell\ge a\right)
\ \text{and}\
\left(\exists M<\infty,\ \exists\ell_1,\ \forall\ell\ge\ell_1:
\min\{t_\ell,d_\ell\}\le M\right).
\]
For every such sequence with \(t_\ell\to\infty\),
\[
\rho(t_\ell,d_\ell)\asymp t_\ell^{-1}
\quad\Longleftrightarrow\quad
\exists M\in\mathbb N,\ \exists\ell_0,\ \forall\ell\ge\ell_0:
d_\ell\le M.
\]

Finally, for every \(0<a\le b<\infty\), there exist constants \(0<c\le C_0<\infty\) and an integer \(D_0\) such that, for every integer \(d\ge\max\{2,D_0\}\) and every \(t>0\) with \(a\le t/d^2\le b\),
\[
c\left[\frac{\log\log(ed)}{\log(ed)}\right]^2
\le \rho(t,d)\le
C_0\left[\frac{\log\log(ed)}{\log(ed)}\right]^2.
\]
The constants and dimension threshold in this last comparison may depend on \(a,b\).
\end{definitionv}

The symmetric law \leanref{sym:G_\theta}{\ensuremath{G_\theta}} in \cref{obj:def:symmetric-law} pairs arm means equally around \(1/2\). The following result states its membership in the causal class and establishes an exact correspondence between private signed observations and private trial records. The correspondence preserves both collection classes and the joint law of the public seed and messages.

\begin{definitionv}{def:symmetric-law}[Symmetric causal law $G_\theta$]
The full-data law \(G_\theta\) is generated by
\[
X\sim\operatorname{Uniform}[d],
\qquad
Y^0\mid X=j\sim\operatorname{Bernoulli}\!\left(\frac{1-\theta_j}{2}\right),
\qquad
Y^1\mid X=j\sim\operatorname{Bernoulli}\!\left(\frac{1+\theta_j}{2}\right),
\]
with \(Y^0\) and \(Y^1\) conditionally independent given \(X\). Independently of \((X,Y^0,Y^1)\), generate the treatment assignment and set the observed outcome by
\[
A\sim\operatorname{Bernoulli}(1/2),
\qquad
Y=Y^{\mathrm{pot}}_A.
\]
\end{definitionv}

\Cref{obj:prop:signed-causal-bridge} identifies the absolute-contrast component with total variation from the uniform reference for every law in the causal class. On the symmetric family, the baseline equals \(1/2\), so variation in welfare is precisely variation in this distance. The assignment bit is fair conditional on the signed input, and the observed outcome is determined by that input and the assignment. Averaging over the assignment therefore reproduces the private transcript of any admissible trial-record protocol on this family. Conversely, a signed-input protocol can be applied to trial records by computing the sign locally. These transformations preserve full-record privacy and history dependence, allowing estimation procedures and decision lower bounds for the signed experiment to inform welfare analysis in both collection classes.

The absolute value also explains the statistical focus on small treatment contrasts. Selecting the better arm contributes \(|\tau_j(P)|/2\) beyond the cell's fair-assignment mean, and this contribution changes slope at a zero contrast. The signed experiment retains the entire contrast vector in its input probabilities, while the target aggregates coordinate magnitudes. Thus the welfare problem combines estimation of the linear baseline with estimation of this nonlinear aggregate.

\subsection{Resource scales and paired procedures}

To state the finite-resource comparisons uniformly across sample sizes, dimensions, and privacy budgets, we use the effective resource \(t=n\varepsilon^2\) and the logarithmic dimension factor \(L=\log(ed)\). The following definition records the sequence conventions used later for the paired minimax values; the paired model and decision criteria are introduced immediately afterward.

\begin{definitionv}{synth_6}[Sequence conclusions for paired minimax values]
For each \(\mathcal C\in\{\mathrm{NI},\mathrm{SI}\}\), let
\[
\mathcal R_{\mathcal C}(n,d,\varepsilon)
=\mathfrak R^{\mathrm{TV}}_{\mathcal C}
(\mathcal P_d^{\mathrm{pair}};n,d,\varepsilon),
\qquad
\mathcal H_{\mathcal C}(n,d,\varepsilon)
=\mathfrak H^{\mathrm{TV}}_{\mathcal C}
(\mathcal P_d^{\mathrm{pair}};n,d,\varepsilon),
\]
where the minimax risk and globally \(0.90\)-honest connected-interval length are defined in \cref{obj:def:tv-private-decision-values}. We say that the sequence-quantified value conclusions hold when the following properties hold for each class.

For every sequence \((n_\ell,d_\ell,\varepsilon_\ell)_{\ell\ge1}\) with
\[
n_\ell,d_\ell\in\{2,3,\ldots\},
\qquad 0<\varepsilon_\ell\le1,
\]
write
\[
t_\ell=n_\ell\varepsilon_\ell^2,\qquad
L_\ell=\log(ed_\ell),\qquad
R_\ell=\mathcal R_{\mathcal C}(n_\ell,d_\ell,\varepsilon_\ell),
\qquad
H_\ell=\mathcal H_{\mathcal C}(n_\ell,d_\ell,\varepsilon_\ell).
\]
Then
\[
R_\ell\longrightarrow0
\quad\Longleftrightarrow\quad
t_\ell\longrightarrow\infty
\ \text{and}\
\frac{\log(1+d_\ell^2/t_\ell)}{L_\ell}\longrightarrow0,
\]
and, separately,
\[
H_\ell\longrightarrow0
\quad\Longleftrightarrow\quad
t_\ell\longrightarrow\infty
\ \text{and}\
\frac{\log(1+d_\ell^2/t_\ell)}{L_\ell}\longrightarrow0.
\]

Define \(R_\ell\asymp t_\ell^{-1}\) to mean that there exist \(0<c\le C_0<\infty\) and \(\ell_0\) such that
\[
c/t_\ell\le R_\ell\le C_0/t_\ell
\qquad(\ell\ge\ell_0).
\]
For every allowed sequence,
\[
R_\ell\asymp t_\ell^{-1}
\quad\Longleftrightarrow\quad
\left(\exists a>0,\ \exists\ell_0,\ \forall\ell\ge\ell_0:
t_\ell\ge a\right)
\ \text{and}\
\left(\exists M<\infty,\ \exists\ell_1,\ \forall\ell\ge\ell_1:
\min\{t_\ell,d_\ell\}\le M\right).
\]
For every allowed sequence satisfying \(t_\ell\to\infty\),
\[
R_\ell\asymp t_\ell^{-1}
\quad\Longleftrightarrow\quad
\exists M\in\mathbb N,\ \exists\ell_0,\ \forall\ell\ge\ell_0:
d_\ell\le M.
\]

Finally, for every \(0<a\le b<\infty\), there exist \(0<c\le C_0<\infty\) and an integer \(D_0\) such that, for every allowed tuple \((n,d,\varepsilon)\) with \(d\ge D_0\) and
\[
a\le\frac{n\varepsilon^2}{d^2}\le b,
\]
both comparisons hold:
\[
c\left[\frac{\log\log(ed)}{\log(ed)}\right]^2
\le\mathcal R_{\mathcal C}(n,d,\varepsilon)
\le C_0\left[\frac{\log\log(ed)}{\log(ed)}\right]^2,
\]
\[
\sqrt c\left|\frac{\log\log(ed)}{\log(ed)}\right|
\le\mathcal H_{\mathcal C}(n,d,\varepsilon)
\le\sqrt{C_0}\left|\frac{\log\log(ed)}{\log(ed)}\right|.
\]
These constants and the dimension threshold may depend on \(a,b\) and the fixed class, and are uniform over the displayed resource domain.
\end{definitionv}

The squared-error scale \leanref{sym:r_{\mathcal C}}{\ensuremath{r_{\mathcal C}(n,d,\varepsilon)}} and length scale \leanref{sym:h_{\mathcal C}}{\ensuremath{h_{\mathcal C}(n,d,\varepsilon)}} in \cref{obj:synth_3} use the same resource function for both collection classes. The defined resource conclusions distinguish a saturated range, an intermediate range, and a dense range, and express consistency and inverse-resource comparisons along arbitrary admissible sequences. They provide the notation for the statistical comparisons made in \cref{obj:thm:uniform-private-value-frontiers,obj:thm:paired-uniform-tv-frontier}.

We also specify the attainment guarantees for the paired experiment. Its input family \leanref{sym:\mathcal P_d^{\mathrm{pair}}}{\ensuremath{\mathcal P_d^{\mathrm{pair}}}} consists of the paired signed laws, and its target \leanref{sym:D_{\mathrm{TV}}}{\ensuremath{D_{\mathrm{TV}}(P_\theta)}} is their total-variation distance from the uniform reference, as defined in \cref{obj:def:tv-corollary-models}. The paired-input protocol class \leanref{sym:\mathfrak K}{\ensuremath{\mathfrak K_{\mathcal C}(n,d,\varepsilon)}} uses the same public randomness, message histories, and one-message privacy requirement on the paired alphabet; its full definition is in \cref{obj:def:tv-private-decision-values}. A paired protocol \leanref{sym:\mathsf K}{\ensuremath{\mathsf K}} generates the private message vector \leanref{sym:\mathbf Z^{\mathsf K}}{\ensuremath{\mathbf Z^{\mathsf K}}}. The estimate \leanref{sym:\widehat D}{\ensuremath{\widehat D}} and interval \leanref{sym:J_{\mathrm{TV}}}{\ensuremath{J_{\mathrm{TV}}}} target the numerical distance. The next definition packages the finite-message attainment properties that the subsequent results establish.

\begin{definitionv}{synth_5}[Resource-selected paired polynomial procedure]
Fix \(n,d\in\{2,3,\ldots\}\), \(0<\varepsilon\le1\), and \(C_0>0\). Put
\[
t=n\varepsilon^2,\qquad L=\log(ed),\qquad
\delta=\tanh(\varepsilon/2),\qquad b=d/\delta,
\qquad m=\lfloor n/3\rfloor.
\]
Use a singleton public seed. For participants \(i=1,\ldots,2m\), release a sign vector \(Z_i\in\{-1,1\}^d\) with conditional probabilities
\[
\mathsf K_i(z\mid(j,s))
=2^{-d}(1+\delta s z_j).
\]
The remaining participants release the constant vector \((-1,\ldots,-1)\). All releases use independent local randomness and are independent of histories.

If \(n=2\), set \(\widehat D=1/8\). If \(n>2\), use the first \(m\) active participants as the pilot block and the next \(m\) as the evaluation block. Define
\[
W^{\mathrm{pil}}_{ij}=bZ_{i,j},
\qquad
W^{\mathrm{ev}}_{ij}=bZ_{m+i,j},
\qquad 1\le i\le m,\quad j\in[d],
\]
and
\[
\sigma^2=\frac{b^2}{m},\qquad
\sigma=\sqrt{\sigma^2},\qquad
H=\log(e+\sigma^2L),\qquad
M_j=\frac1m\sum_{i=1}^mW^{\mathrm{pil}}_{ij},
\qquad
\overline W_j=\frac1m\sum_{i=1}^mW^{\mathrm{ev}}_{ij}.
\]
For the evaluation block, define the distinct-person moments
\[
\mathsf U_{0j}=1,\qquad
\mathsf U_{vj}
=\binom mv^{-1}
\sum_{\substack{J\subseteq\{1,\ldots,m\}\\ |J|=v}}
\prod_{i\in J}W^{\mathrm{ev}}_{ij},
\qquad 1\le v\le m.
\]
Their finite computation uses elementary-symmetric recursions.

For each even degree \(D\ge2\) selected below, define
\[
\mathcal T_0(x)=1,\qquad
\mathcal T_1(x)=x,\qquad
\mathcal T_{v+1}(x)=2x\mathcal T_v(x)-\mathcal T_{v-1}(x)
\quad(v\ge1),
\]
\[
p_D^{\mathrm{abs}}(x)
=\frac2\pi+\frac4\pi
\sum_{v=1}^{D/2}
\frac{(-1)^{v+1}\mathcal T_{2v}(x)}{4v^2-1}
=\sum_{v=0}^Dc_{D,v}x^v,
\]
and, for a positive radius \(a_{\mathrm{rad}}\),
\[
\widehat p_{a_{\mathrm{rad}},D,j}
=\sum_{v=0}^D
c_{D,v}a_{\mathrm{rad}}^{\,1-v}\mathsf U_{vj}.
\]

For \(n>2\), define \(\widehat F\) by the following ordered branches, taking the first applicable branch:
\[
\widehat F=
\begin{cases}
\displaystyle\frac1d\sum_{j=1}^d|\overline W_j|,
&L<4096,\\[6pt]
\displaystyle\frac1d\sum_{j=1}^d
\left[
\widehat p_{a_{\mathrm{rad}},D,j}
\mathbf1_{\{|M_j|\le T_*\}}
+\operatorname{sign}(M_j)\overline W_j
\mathbf1_{\{|M_j|>T_*\}}
\right],
&L\ge4096,\ t\ge d^2L,\\[6pt]
1/4,
&L\ge4096,\ t<d^2L,\ L/(1024H)<4,\\[4pt]
\displaystyle\frac1d\sum_{j=1}^d
\widehat p_{a_{\mathrm{rad}},D,j},
&L\ge4096,\ t<d^2L,\ L/(1024H)\ge4.
\end{cases}
\]
In the second branch, set
\[
D=2\left\lfloor\frac{L}{2048}\right\rfloor,
\qquad T_*=4\sigma\sqrt L,\qquad
a_{\mathrm{rad}}=8\sigma\sqrt L,
\qquad \operatorname{sign}(0)=0.
\]
In the fourth branch, set
\[
a_{\mathrm{rad}}=\frac12,
\qquad D=2\left\lfloor\frac{L}{2048H}\right\rfloor.
\]
For \(n>2\), set
\[
\widehat D=\max\left\{0,\min\left\{\frac14,\frac{\widehat F}{2}\right\}\right\}.
\]
For every \(n\ge2\), report
\[
q=\sqrt{10C_0\rho(t,d)},\qquad
J_{\mathrm{TV}}
=\left[
\max\{0,\widehat D-q\},
\min\{1/4,\widehat D+q\}
\right].
\]

We say that the concrete procedure guarantees hold with constant \(C_0\) when this fixed protocol belongs to \(\mathfrak K_{\mathrm{NI}}(n,d,\varepsilon)\) and, for every \(\theta\in[-1/2,1/2]^d\),
\[
\begin{aligned}
\mathbb E_{P_\theta,\mathsf K}
\left[(\widehat D-D_{\mathrm{TV}}(P_\theta))^2\right]
&\le C_0\rho(t,d),\\
\Pr_{P_\theta,\mathsf K}
\{D_{\mathrm{TV}}(P_\theta)\in J_{\mathrm{TV}}\}
&\ge0.90,\\
\mathbb E_{P_\theta,\mathsf K}
\left[\operatorname{length}J_{\mathrm{TV}}\right]
&\le C_0\sqrt{\rho(t,d)}.
\end{aligned}
\]
Here \(P_\theta(j,s)=(1+s\theta_j)/(2d)\) and
\(D_{\mathrm{TV}}(P_\theta)=(2d)^{-1}\sum_j|\theta_j|\).
Both decisions are Borel functions of the transcript, and the reported interval is nonempty, connected, closed, and contained in \([0,1/4]\).
\end{definitionv}

The attainment properties in \cref{obj:synth_4} require a finite sign-vector message from each participant and place both decisions on the paired target's natural range. They combine squared-error control, global coverage, and expected connected-interval length under a common protocol. A more explicit specification selects its polynomial estimate from public resource parameters. In that construction, the privacy attenuation \leanref{sym:\delta}{\ensuremath{\delta}} determines the scaled-message magnitude \leanref{sym:b}{\ensuremath{b}}, and the block size \leanref{sym:m}{\ensuremath{m}} separates pilot information from evaluation information.

\begin{definitionv}{synth_4}[Finite sign-vector attainment]
Fix \(n,d\in\{2,3,\ldots\}\), \(0<\varepsilon\le1\), and \(C_0>0\). Let
\[
P_\theta(j,s)=\frac{1+s\theta_j}{2d},
\qquad
\theta\in[-1/2,1/2]^d,\quad j\in[d],\quad s\in\{-1,1\}.
\]
The paired family and its total-variation target are
\[
\mathcal P_d^{\mathrm{pair}}
=\{P_\theta:\theta\in[-1/2,1/2]^d\},
\qquad
D_{\mathrm{TV}}(P_\theta)
=\frac1{2d}\sum_{j=1}^d|\theta_j|.
\]
We say that the general attainment guarantees hold with constant \(C_0\) when there exist a protocol
\[
\mathsf K\in\mathfrak K_{\mathrm{NI}}(n,d,\varepsilon),
\]
a Borel estimator \(\widehat D\), and a Borel interval decision \(J_{\mathrm{TV}}\) with the following properties.

Each participant's message space admits a measurable bijection with \(\{-1,1\}^d\). The protocol is a one-person-one-message pure-local protocol on the original alphabet \(\mathcal A_d=[d]\times\{-1,1\}\), as specified in \cref{obj:def:tv-private-decision-values}. In particular, its kernels are constant in the history and satisfy
\[
\mathsf K_i(E\mid(j,s),\eta,r)
\le e^\varepsilon
\mathsf K_i(E\mid(j',s'),\eta,r)
\]
for every participant, public seed, history, Borel message event, and pair of original inputs.

The estimator depends on the transcript alone: there is a Borel function \(f\) such that
\[
\widehat D(\mathbf Z^{\mathsf K},R,U)=f(\mathbf Z^{\mathsf K}).
\]
The interval \(J_{\mathrm{TV}}(\mathbf Z^{\mathsf K},R,U)\) is nonempty, connected and closed, has ordered Borel endpoints, and is contained in \([0,1/4]\).

For every \(\theta\in[-1/2,1/2]^d\), under \(n\) iid inputs from \(P_\theta\),
\[
\begin{aligned}
\mathbb E_{P_\theta,\mathsf K}
\left[(\widehat D-D_{\mathrm{TV}}(P_\theta))^2\right]
&\le C_0\,\rho(n\varepsilon^2,d),\\
\Pr_{P_\theta,\mathsf K}
\{D_{\mathrm{TV}}(P_\theta)\in J_{\mathrm{TV}}\}
&\ge0.90,\\
\mathbb E_{P_\theta,\mathsf K}
\left[\operatorname{length}J_{\mathrm{TV}}\right]
&\le C_0\sqrt{\rho(n\varepsilon^2,d)}.
\end{aligned}
\]
Probabilities and expectations include input sampling, channel randomness, the independent public seed, and independent analyst randomization. The same protocol belongs to \(\mathfrak K_{\mathrm{SI}}(n,d,\varepsilon)\).
\end{definitionv}

The construction in \cref{obj:synth_5} fixes collection independently of histories and uses resource-dependent branches for the statistical decision. Its scaled messages \leanref{sym:W}{\ensuremath{W^{\mathrm{pil}}}} and \(W^{\mathrm{ev}}\) support pilot averages and distinct-person moments \leanref{sym:\mathsf U}{\ensuremath{\mathsf U_{vj}}}. In the dense, large-dimension branch, the pilot selects between a polynomial estimate near zero and a signed evaluation mean away from zero. The other branches use an absolute empirical mean, a constant, or a polynomial across the full contrast range, according to the displayed resource conditions. This deliberately conservative construction serves as an explicit uniform rate-attainment witness. The concrete guarantees are the three uniform decision bounds specified in the definition; their establishment is part of \cref{obj:thm:paired-uniform-tv-frontier}.

Finally, the paired minimax squared-error criterion \leanref{sym:\mathfrak R^{\mathrm{TV}}}{\ensuremath{\mathfrak R^{\mathrm{TV}}_{\mathcal C}}} and honest-length criterion \leanref{sym:\mathfrak H^{\mathrm{TV}}}{\ensuremath{\mathfrak H^{\mathrm{TV}}_{\mathcal C}}} are defined in \cref{obj:def:tv-private-decision-values}. The next definition records the sequence comparisons for these values that accompany the finite-resource bounds. It keeps the quantifiers explicit when dimension and privacy vary with sample size.

\begin{propositionv}{prop:signed-causal-bridge}[Signed causal experiment equivalence]
Fix a sample size \(n\in\mathbb N\), a stratum count \(d\in\mathbb N\), and a sampling scheme satisfying:
\begin{itemize}
\item \textbf{(Dimension.)} \(d\ge 2\).
\item \textbf{(Independent participants.)} \cref{obj:ass:iid-people}.
\item \textbf{(Independent randomization.)} \cref{obj:ass:independent-randomness}.
\end{itemize}
Let \(\mathcal V_d\) be the class of causal probability laws in \cref{obj:def:causal-model}. For the stratum \(X\), binary assignment \(A\), observed outcome \(Y\), and potential outcomes \(Y^0,Y^1\), write
\[
S=(2A-1)(2Y-1),\qquad
\mu_{aj}(P)=\mathbb E_P[Y^a\mid X=j],\qquad
\tau_j(P)=\mu_{1j}(P)-\mu_{0j}(P).
\]
Write \(\tau(P)\) for the vector of these contrasts, and define the baseline mean, first-best value, and average absolute contrast by
\[
B(P)=\mathbb E_P[Y],\qquad
V(P)=\frac1d\sum_{j=1}^d\max\{\mu_{0j}(P),\mu_{1j}(P)\},
\qquad
F(\theta)=\frac1d\sum_{j=1}^d|\theta_j|.
\]
The paired signed law \(P_\theta\) has masses
\[
P_\theta(j,s)=\frac{1+s\theta_j}{2d},
\qquad j\in[d],\quad s\in\{-1,1\},
\]
and \(\mathsf U_{2d}\) is the uniform law on this paired alphabet.

For every \(P\in\mathcal V_d\) and every \(j\in[d]\),
\[
\mathbb E_P[S\mid X=j]=\tau_j(P),
\qquad
\mathcal L_P(X,S)=P_{\tau(P)},
\]
and
\[
V(P)
=B(P)+\frac{F(\tau(P))}{2}
=B(P)+\operatorname{TV}(P_{\tau(P)},\mathsf U_{2d}).
\]

For every \(\theta\in[-1/2,1/2]^d\), the symmetric law \(G_\theta\) of \cref{obj:def:symmetric-law} is a probability law in \(\mathcal V_d\), and
\[
V(G_\theta)=\frac12+\frac{F(\theta)}2.
\]
Moreover, for every \(j\in[d]\), \(s\in\{-1,1\}\), and \(a\in\{0,1\}\),
\[
G_\theta(X=j,S=s,A=a)=\frac{P_\theta(j,s)}2,
\]
and
\[
2Y-1=S(2A-1)
\qquad G_\theta\text{-almost surely}.
\]

For every privacy parameter \(\varepsilon\in\mathbb R\), the following two protocol transformations preserve both the sequential and noninteractive classes of \cref{obj:def:sequential-class,obj:def:noninteractive-class}. Let \(\mathcal H_i\) denote the message-history space preceding participant \(i\), and let \(\mathcal R\) denote the declared space of the public-seed random variable \(R\). For every original-record protocol \(Q=(Q_i)_{i=1}^n\), define its paired-alphabet protocol \(\mathsf K\) by averaging over the fair assignment bit:
\[
\mathsf K_i(\,\cdot\mid(j,s),h,r)
=
\frac12\sum_{a\in\{0,1\}}
Q_i\!\left(
\,\cdot\mid
\left(j,a,\frac{1+s(2a-1)}2\right),h,r
\right),
\qquad
h\in\mathcal H_i,\quad r\in\mathcal R.
\]
Sequential membership of \(Q\) at \(\varepsilon\) implies sequential membership of \(\mathsf K\) at \(\varepsilon\), and noninteractive membership of \(Q\) implies noninteractive membership of \(\mathsf K\) at \(\varepsilon\). For every \(\theta\in[-1/2,1/2]^d\), the joint law of the public seed \(R\) and message vector \(\mathbf Z\) generated by \(Q\) under \(G_\theta\) and the fixed sampling scheme equals the joint law generated by \(\mathsf K\) from \(n\) independent \(P_\theta\) inputs and an independent public seed with the same law.

Conversely, for every paired-alphabet protocol \(\mathsf K\), define its original-record protocol \(Q\) by
\[
Q_i(\,\cdot\mid(j,a,y),h,r)
=
\mathsf K_i(\,\cdot\mid(j,(2a-1)(2y-1)),h,r),
\qquad
h\in\mathcal H_i,\quad r\in\mathcal R.
\]
Sequential membership of \(\mathsf K\) at \(\varepsilon\) implies sequential membership of \(Q\) at \(\varepsilon\), and noninteractive membership of \(\mathsf K\) implies noninteractive membership of \(Q\) at \(\varepsilon\). For every \(\theta\in[-1/2,1/2]^d\), the joint public-seed and message law generated by this \(Q\) under \(G_\theta\) and the fixed sampling scheme equals that generated by \(\mathsf K\) from \(n\) independent \(P_\theta\) inputs and an independent public seed with the same law.
\end{propositionv}

The comparisons collected in \cref{obj:synth_6} distinguish vanishing estimation error, shrinking honest intervals, and inverse-resource squared error. Their uniform resource-domain formulation also specifies the dimension-dependent scale when effective resource is proportional to \(d^2\). Together with the exact signed-experiment correspondence in \cref{obj:prop:signed-causal-bridge}, these definitions provide the population, observation, and decision framework for the welfare and total-variation results.
